In this work, we first explored the state of the art regarding the modeling of energy communities, the gain allocation methods and the fairness problematics. A lot of different models of energy communities have been proposed in the literature in order to respond to the different technical problematics. The sociological aspect of the energy communities is a very important aspect to take into account as the success of the communities depends on the involvement of the people in their development. Threw sociological profiles the models have tried to take this aspect into consideration directly in the operation of the community. Furthermore, the way the gain is allocated has very important implications on the fairness within a community. 

However, in our view, the literature often lacks clarity regarding the significance of the social aspect. This study, therefore, aimed to highlight the impact of sociological profiles on community operations and to illustrate how different gain allocation methods influence fairness within the community.

For this work, we chose to make an operation model of a REC that could be easily optimized in a fully centralized way and in a fully decentralized way. This way we can compute the gain of the community and allocate it using different allocation functions. The model is built from the bottom up, we first model the different devices and flexibility services that can be found in a house. Then we aggregate these models to make the model of a member of the community. These models can then be aggregated again to make the model of the community. 

The gains can then be computed threw several methods. In the state of the art, a lot of them are presented, for simplicity's sake, only global allocation functions were implemented in this study. The three different functions used are : the Shapley value, a proportional allocation and an equal allocation. These functions have been chosen in order to represent different properties in terms of fairness. The Shapley value is a "meritocratic" allocation function, it gives more gain to the members depending on their marginal contribution. The proportional allocation respects the Nash equilibrium, it gives to each member a gain proportional to its contribution in reducing the global cost of the community. This should respect the "min-max" fairness as defined in the state of the art. Finally, the equal allocation gives the same gain to all the members of the community, it is an "egalitarian" allocation function.

In order to provide some visuals, we developed an example of community with 5 members by creating a synthetic dataset. To do so, first we developed some simplified models of different objects in a house. By coupling these with randomly generated presence profiles, we were able to generate random members for a community. After building this community, we could finally compute different values.

Firstly, we presented the Pareto front of the model regarding the different sociological profiles. It shows that there is a clear trade-off between the economic/ecological objective function and the comfort objective function. And the economic and ecological objective functions are very correlated. This was expected and would be more interesting if we were to take into account the investment of the different devices.

Then we presented the gain allocation of the different functions. For the equality function, we could have understood easily what it would imply. Now, looking at the Shapley value, we could see some interesting results. It is meritocratic, as it gives more gain to the members that contribute more to the gain of the community. But we can see that there is a penalty towards the members that have a higher consumption as, even if they have a high production rate and a battery, their allocated gain is lower than one of the others without production. The proportional allocation reaches a result that looks hard to interpret. Knowing that it should reach a Nash equilibrium, we can understand better the results. As such it can be seen as a "max-min" fairness allocation function.

The results shown shows how different allocation functions can lead to very different gain allocation as they follow a different idea of fairness. In real case application, this choice of the allocation methods should be made carefully, for example, by involving the community in the decision. It should be important for the members to understand the implications of the different methods in order to choose the one that best fit their idea of fairness. This could improve their engagement in the community and thus the success of the community.

This study could be completed by implementing different fairness indexes that are presented in \cite{article_1, these_new, article_8, article_2} in order to evaluate the different allocation methods. It was not done here, as the focus of the work was more on visualizing the effect of the allocation methods. But when comparing more allocation methods, these indexes could be very useful, even if it is not always easy to interpret what they are actually measuring. Also, now that we have this work on 3 different allocation methods, it could be interesting to compare more widely the different methods and their properties in terms of fairness. Finally, different sociological profiles are used in the community. In this work, a resulting sociological profile is computed as the average of the sociological profiles of the members, and so maybe the final sociological profile of the community does not fit the different members. It could be interesting to look into the aspect of fairness regarding this. For example, some none fully centralized models could take into account the different sociological profiles of the members in order to find a compromise between the different profiles. 

Another aspect that was not looked at in this work is the influence of the number of people in the community and the time complexity of the operation. For example, the Shapley value has an exponential complexity regarding the number of members in the community, and is often deemed impossible to compute for more than 20 members. In general, the computation time was not looked into in this work, but it can become a limitation when increasing the number of members within the community.

\subsection{Lien avec les cours de ERI}

Ce travail fait référence aux cours sur les communautés énergétiques locales dans la première partie de l'année ainsi qu'aux cours d'optimisation que nous avons eus. Un modèle d'optimisation plus complexe a été développé que les modèles simples que nous avons utilisés en classe. Par conséquent, j'ai pu explorer davantage le développement de modèle par rapport à ce qui a été vu en classe. De plus, ce travail essaie de considérer d'autres aspects que les problématiques techniques que nous avons vues en cours concernant le réseau électrique. Ainsi, il a été très intéressant d'explorer des aspects sociaux tels que les problématiques d'équité.

\subsection{Analyse reflexive}

Avec ce travail, j'ai beaucoup appris sur la modélisation des communautés énergétiques, et les différentes problématiques autour de celle-ci, qu'elles soient techniques ou sociales. D'un point de vue pratique, j'ai développé beaucoup mes capactités de lecture et de compréhension d'article scientifique. Et surtout, je pense avoir fait des progrès dans la restitution de mon travail au travers du rapport bibliographique et de ce rapport final. Et ce notamment en anglais.

J'ai aussi continué de développer mes compétences en modélisation et optimisation. Mais aussi en code informatique, car je pense être arrivé à produire un code plus propre et exploitable par rapport à d'autres de mes projets précédents. 

Je reconnais cependant que je peux encore faire des progrès sur l'organisation du temps de travail. J'ai eu du mal à répartir mon travail dans le temps plutôt que de travailler par gros blocs courts.